% from ega4-2-fr.tex
\oldpage[IV-2]{6}
\begin{proposition}[2.1.3]
\label{IV.2.1.3-fr}
Soient $f:X\to Y$ un morphisme de préschémas, $\mathscr F$ un
$\mathcal O_X$-Module quasi-cohérent. Les conditions suivantes sont
équivalentes~:
\begin{enumerate}
\item[\emph{a)}] Pour tout changement de base $g:Y'\to Y$, si l'on pose
$X'=X\times_Y Y'$, le foncteur
$\mathscr G'\mapsto\mathscr F\otimes_Y\mathscr G'$ de la catégorie des
$\mathcal O_{Y'}$-Modules quasi-cohérents dans celle des
$\mathcal O_{X'}$-Modules quasi-cohérents, est exact.
\item[\emph{a$'$)}] La condition \emph{a)} est vérifiée pour tous les
morphismes canoniques $g:\operatorname{Spec}(\mathcal O_y)\to Y$ (I, 2.4.1),
où $y$ parcourt $Y$.
\item[\emph{b)}] $\mathscr F$ est $f$-plat.
\end{enumerate}
\end{proposition}

Les questions étant locales sur $X$ et $Y$, on peut se borner au cas où
$Y=\operatorname{Spec}(B)$, $X=\operatorname{Spec}(A)$,
$\mathscr F=\widetilde M$, où $M$ est un $A$-module. Il est clair que
\emph{a)} entraîne \emph{a$'$)}~; la condition \emph{a$'$)} entraîne que
pour tout $x\in X$, le foncteur $N\mapsto M_{\mathfrak n}\otimes_{B_{\mathfrak n}}N$
est exact dans la catégorie des $B_{\mathfrak n}$-modules, $\mathfrak n$
étant l'idéal $\mathfrak j_{f(x)}$ de $B$~; cela signifie que
$M_{\mathfrak n}$ est un $B_{\mathfrak n}$-module plat, et il résulte de
$(0_{\mathrm I},6.3.3)$ et de (2.1.2) que $\mathscr F$ est $f$-plat. Enfin,
pour voir que \emph{b)} entraîne \emph{a)}, on peut aussi se borner au cas où
$Y'=\operatorname{Spec}(A')$ est affine et $\mathscr G'=\widetilde{N'}$, où
$N'$ est un $A'$-module~; la conclusion provient alors encore de (2.1.2) et
de la définition de la platitude, puisque
$(M\otimes_A N')^{\sim}=\mathscr F\otimes_Y\mathscr G'$.

\begin{proposition}[2.1.4]
\label{IV.2.1.4-fr}
Soient $f:X\to Y$, $g:Y'\to Y$ deux morphismes de préschémas, $\mathscr F$
un $\mathcal O_X$-Module quasi-cohérent~; posons $X'=X\times_Y Y'$,
$f'=f_{(Y')}:X'\to Y'$, $\mathscr F'=\mathscr F\otimes_Y\mathcal O_{Y'}$,
et soit $g'$ la projection canonique $X'\to X$. Soient $x'$ un point de
$X'$, $x=g'(x')$, $y'=f'(x')$, $y=g(y')=f(x)$. Si $\mathscr F$ est $f$-plat
au point $x$, $\mathscr F'$ est $f'$-plat au point $x'$~; en particulier si
$\mathscr F$ est $f$-plat, $\mathscr F'$ est $f'$-plat~; si $f$ est plat,
$f'$ est plat.
\end{proposition}

Il suffit de prouver la première assertion~; appliquant trois fois (I, 3.6.5),
ainsi que (I, 2.4.4), on peut se ramener au cas où
$Y=\operatorname{Spec}(\mathcal O_y)$, $X=\operatorname{Spec}(\mathcal O_x)$,
$Y'=\operatorname{Spec}(\mathcal O_{y'})$, $\mathscr F=\widetilde M$, où
$M=\mathscr F_x$~; l'hypothèse et (2.1.2) entraînent alors que $\mathscr F$
est $f$-plat, autrement dit, on est ramené à prouver un cas particulier de la
seconde assertion, et cette dernière découle aussitôt de (2.1.3).

\begin{proposition}[2.1.5]
\label{IV.2.1.5-fr}
Considérons un diagramme commutatif de morphismes de préschémas
\[
\xymatrix{
  X \ar[d]_f & X' \ar[l]_{g'} \ar[d]^{f'} \\
  Y & Y' \ar[l]^g \ar[d]^h \\
  & Z
}
\]
où $X'=X\times_Y Y'$ et $f'=f_{(Y')}$. Soit $x'$ un point de $X'$, et
posons $x=g'(x')$, $y'=f'(x')$, $y=f(x)=g(y')$, $z=h(y')$. Soient
$\mathscr F$ un $\mathcal O_X$-Module quasi-cohérent $f$-plat au point $x$
(resp. $f$-plat), et $\mathscr G'$ un $\mathcal O_{Y'}$-Module quasi-cohérent
$h$-plat au point $y'$ (resp. $h$-plat)~; alors
$\mathscr F\otimes_{Y'}\mathscr G'$ est un $\mathcal O_{X'}$-Module
quasi-cohérent $(hf')$-plat au point $x'$ (resp. $(hf')$-plat).
\end{proposition}

Comme dans (2.1.4), on se ramène au cas où
$X=\operatorname{Spec}(\mathcal O_x)$, $Y=\operatorname{Spec}(\mathcal O_y)$,
